% =============================================================================
% Presentación de Defensa de Tesis
% Aumentación de Datos con LLMs: Filtrado Geométrico para Clasificación de Texto
% Benjamín Schindler — Universidad Adolfo Ibáñez — Abril 2026
% Profesor guía: Gonzalo Ruz
% Compilar con: xelatex (preferido) o lualatex
% =============================================================================

\documentclass[10pt,aspectratio=169]{beamer}

% --- Tema Metropolis con paleta UAI ---
\usetheme[
  progressbar=frametitle,
  sectionpage=progressbar,
  numbering=fraction,
  block=fill
]{metropolis}

% --- Paquetes base ---
\usepackage[spanish,es-tabla,es-noindentfirst]{babel}
\usepackage{fontspec}
\usepackage{amsmath,amssymb}
\usepackage{booktabs,multirow,multicol}
\usepackage{tikz}
\usetikzlibrary{shapes,arrows.meta,positioning,fit,calc,backgrounds,shadows.blur}
\usepackage{graphicx}
\usepackage{xcolor}
\usepackage{tcolorbox}
\usepackage{appendixnumberbeamer}
\usepackage{pifont}
\usepackage{caption}
\captionsetup{font=scriptsize, labelfont={bf,scriptsize}}

% --- Fuente: Cantarell (similar a Fira Sans, instalada en sistema) ---
\setsansfont{Cantarell}[
  Extension       = .otf,
  UprightFont     = *-Regular,
  ItalicFont      = *-Regular,
  BoldFont        = *-Bold,
  BoldItalicFont  = *-Bold,
  Path            = /usr/share/fonts/opentype/cantarell/,
  Scale           = 0.95
]
\setmonofont{DejaVu Sans Mono}[Scale=0.85]

% --- Estilos custom UAI ---
\usepackage{style/uai-colors}
\usepackage{style/beamer-blocks}

% --- Configuración Metropolis adicional ---
\setbeamertemplate{section in toc}[sections numbered]
\metroset{titleformat=regular}

% --- Section page custom: título grande + indicador lateral ---
\setbeamertemplate{section page}{%
  \begin{centering}
  \vspace{0.3em}
  \begin{minipage}{0.94\textwidth}
    \begin{columns}[c]
      \begin{column}{0.52\textwidth}
        \raggedright
        {\color{uaiorange}\fontsize{42}{46}\selectfont\bfseries\thesecnum.}\\[0.3em]
        {\color{uaibluedark}\Large\bfseries\insertsectionhead}
      \end{column}
      \begin{column}{0.46\textwidth}
        \sectionnav{\value{secnum}}
      \end{column}
    \end{columns}
  \end{minipage}
  \end{centering}%
}
\AtBeginSection[]{%
  \stepcounter{secnum}%
  \begin{frame}[plain,c,noframenumbering]\sectionpage\end{frame}%
}

% --- Ruta de figuras ---
\graphicspath{{figuras/}{logos/}}

% --- Comandos abreviados ---
\newcommand{\dataset}[1]{\texttt{\small #1}}
\newcommand{\filter}[1]{\textit{#1}}
\newcommand{\good}{\textcolor{uaigreen}{\ding{51}}}
\newcommand{\bad}{\textcolor{uaiorange}{\ding{55}}}

% --- Información título ---
\title{Aumentación de Datos con LLMs}
\subtitle{Filtrado Geométrico para Clasificación de Texto en Pocos Ejemplos}
\author{\textbf{Benjamín Schindler}}
\institute{
  Magíster en Ciencia de Datos\\
  Facultad de Ingeniería y Ciencias\\
  Universidad Adolfo Ibáñez\\[4pt]
  \footnotesize Profesor guía: Gonzalo Ruz
}
\date{Abril 2026}
\titlegraphic{\hfill
\includegraphics[height=1.2cm]{logos/fic.png}}

% =============================================================================
\begin{document}

% --- Bloque 1: Apertura ---
% =============================================================================
% Bloque 1: Apertura (slides 1–2)
% =============================================================================

% --- Slide 1: Portada ---
\maketitle

% --- Slide 2: Roadmap (lista con descripciones, sin tiempos) ---
\begin{frame}{Hoja de ruta}
\small
\setbeamercolor{enumerate item}{fg=uaiorange}
\begin{enumerate}\setlength\itemsep{4pt}
  \item \textbf{Introducción} — el régimen \emph{few-shot} y un escenario representativo donde este trabajo ayuda.

  \item \textbf{Problema} — escasez de datos, el marco clásico del re-muestreo y la pregunta central.

  \item \textbf{Estado del arte} — tres familias de aumentación, tres deficiencias persistentes y una década de evolución.

  \item \textbf{Hipótesis y objetivos} — qué afirmamos y cómo se va a verificar.

  \item \textbf{Metodología} — cinco filtros geométricos, ponderación suave y un pipeline desacoplado del LLM.

  \item \textbf{Resultados} — \alert{8.000+ configuraciones}, +2.25 pp sobre SMOTE y un hallazgo contra-intuitivo.

  \item \textbf{Discusión} — utilidad práctica, limitaciones honestas y trabajo futuro.
\end{enumerate}
\end{frame}


% --- Bloque 2: Introducción ---
% =============================================================================
% Bloque 2: Introducción (slides 3–4)
% =============================================================================

\section{Introducción}

% --- Slide 3: Hook + definición + caso real (sin spoiler) ---
\begin{frame}{¿Cuántas muestras se necesitan para clasificar texto?}
\begin{columns}[T]

\begin{column}{0.55\textwidth}
\small
Imagina entrenar un clasificador de \textbf{6 emociones} con apenas\\
\textbf{10 ejemplos por clase}.
\begin{itemize}\setlength\itemsep{3pt}
  \item Total: \alert{60 muestras} en un espacio de \alert{768 dimensiones}.
  \item SMOTE: interpola vectores; no produce texto coherente.
  \item Un LLM genera miles de oraciones plausibles\dots\\
        \cita{pero, ¿son útiles todas?}
\end{itemize}

\vspace{0.5em}
\begin{insight}
\scriptsize
Aprender fronteras de decisión con $60$ puntos en $\mathbb{R}^{768}$ es estadísticamente débil. La aumentación de datos no es opcional: es \textbf{necesaria}.
\end{insight}
\end{column}

\begin{column}{0.45\textwidth}
\setbeamercolor{block title}{fg=white,bg=uaibluedark}
\begin{block}{\small Few-shot text classification}
\scriptsize
Aprender a clasificar texto con un número reducido de ejemplos etiquetados por clase, típicamente $k \leq 50$.\\[3pt]
\textbf{Esta tesis estudia $k \in \{10, 25, 50\}$.}
\end{block}

\vspace{0.6em}
\setbeamercolor{block title}{fg=white,bg=uaiorange}
\begin{block}{\small Caso real: \dataset{emotion}}
\scriptsize
6 clases: alegría, tristeza, miedo, ira, amor, sorpresa.\\
Entrenamiento: \alert{\textbf{60 mensajes}}. Evaluación: $\sim$1.000.\\[2pt]
\textit{La situación del ejemplo es real y se mide en esta tesis.}
\end{block}

\vspace{0.5em}
\centering
{\scriptsize\itshape ¿Cómo aumentar sin etiquetar miles más?}
\end{column}

\end{columns}
\end{frame}

% --- Slide 4: Escena situacional (donde esta tesis es útil) ---
\begin{frame}{Pocos datos, apuro y cero presupuesto: una escena habitual}
\begin{columns}[T]

\begin{column}{0.55\textwidth}
\small
\textbf{Lunes, 9:00.} Tu jefe te pide un clasificador en producción \textbf{el viernes}.

\vspace{0.5em}
\textbf{Lo que tienes:}
\begin{itemize}\setlength\itemsep{2pt}
  \item 12 categorías de mensajes de pacientes (salud digital).
  \item $\sim$15 mensajes etiquetados por categoría desde el CRM.
  \item Cero presupuesto para etiquetado externo.
  \item Un experto clínico ocupado hasta fin de mes.
\end{itemize}

\vspace{0.4em}
\textbf{Lo que el cliente exige:}
\begin{itemize}\setlength\itemsep{2pt}
  \item Macro-F1 $\geq 75\%$ sobre las 12 categorías.
  \item Demo funcionando en \textbf{4 días}.
\end{itemize}

\vspace{0.4em}
\cita{La escena se repite: spam emergente, soporte en idiomas nuevos, triaje clínico\dots}
\end{column}

\begin{column}{0.45\textwidth}
\centering
\vspace{0.4em}
\begin{tikzpicture}[every node/.style={font=\scriptsize, align=center}]
  \node[draw=uaibluedark, fill=white, rounded corners, very thick,
        text width=3.6cm, minimum height=1.1cm] at (0, 3.2)
        {\textbf{Disponible}\\\scriptsize 12 clases $\times$ 15 ej.\\$= 180$ muestras};

  \draw[->, very thick, uaigray] (0, 2.55) -- (0, 2.05);

  \node[draw=uaiorange, fill=uaiorange!10, rounded corners, very thick,
        text width=3.6cm, minimum height=1.1cm] at (0, 1.3)
        {\textbf{Necesario}\\\scriptsize macro-F1 $\geq 75\%$\\en 4 días};

  \node[font=\Large\bfseries, color=uaiorange!85!black] at (0, 0.0) {?};
  \node[font=\small\bfseries, color=uaiorange!85!black] at (0, -0.5) {¿Cómo cerrar la brecha?};
\end{tikzpicture}
\end{column}

\end{columns}

\mensaje{Esta tesis ofrece una respuesta concreta para escenarios como este.}
\end{frame}


% --- Bloque 3: Problema ---
% =============================================================================
% Bloque 3: Problema (Escasez + Re-muestreo + Deficiencias + Pregunta)
% =============================================================================

\section{Problema}

% --- Slide: Escasez de datos en la práctica (movida desde intro) ---
\begin{frame}{Escasez de datos: un problema cotidiano}
\begin{columns}[T]

\begin{column}{0.5\textwidth}
\small
\textbf{Casos reales que sufren pocos ejemplos:}
\begin{itemize}
  \item Detección de \textbf{discurso de odio} en idiomas minorizados.
  \item Triaje clínico, clasificación de notas médicas.
  \item Categorización de \textbf{tickets de soporte} en empresas.
  \item Spam emergente y nuevos esquemas de fraude.
\end{itemize}
\vspace{0.4em}
\textbf{Por qué duele:}
\begin{itemize}
  \item Etiquetado humano = \alert{caro, lento y sesgado}.
  \item Multiclase $+$ desbalance $\Rightarrow$ macro F1 colapsa.
  \cita{Gopali 2024; Whitehouse 2023}
\end{itemize}
\end{column}

\begin{column}{0.5\textwidth}
\centering
\begin{tikzpicture}[every node/.style={font=\scriptsize}]
  % Iconos representando dominios
  \node[draw=uaiblue, rounded corners, fill=uaibluelight!15, minimum width=2.3cm, minimum height=1cm, text width=2.1cm, align=center] (a) at (-1.3,3) {\textbf{Hate speech}\\\tiny minoría: $<$30 ej.};
  \node[draw=uaiblue, rounded corners, fill=uaibluelight!15, minimum width=2.3cm, minimum height=1cm, text width=2.1cm, align=center] (b) at (1.3,3) {\textbf{Clínico}\\\tiny rare disease};
  \node[draw=uaiblue, rounded corners, fill=uaibluelight!15, minimum width=2.3cm, minimum height=1cm, text width=2.1cm, align=center] (c) at (-1.3,1.5) {\textbf{Soporte TI}\\\tiny producto nuevo};
  \node[draw=uaiblue, rounded corners, fill=uaibluelight!15, minimum width=2.3cm, minimum height=1cm, text width=2.1cm, align=center] (d) at (1.3,1.5) {\textbf{Spam}\\\tiny esquema emergente};
  \node[draw=uaiorange, very thick, rounded corners, fill=uaiorange!15, minimum width=4.8cm, minimum height=0.8cm, text width=4.6cm, align=center] at (0,0.2) {\textbf{Común:} 10--50 ejemplos por clase};
\end{tikzpicture}
\end{column}

\end{columns}
\end{frame}

% --- Slide: Re-muestreo (movida desde Estado del Arte) ---
\begin{frame}{Re-muestreo: el paraguas que une a SMOTE y a los LLMs}
\small
La estadística clásica enmarca la escasez desde hace décadas como un problema de \emph{re-muestreo}. Ambas familias caben aquí.

\vspace{0.5em}
\begin{columns}[T]

\begin{column}{0.50\textwidth}
\setbeamercolor{block title}{fg=white,bg=uaiblue}
\begin{block}{Undersampling}
\scriptsize
\textbf{Reduce} la clase mayoritaria.\\[2pt]
\textbf{Ejemplos:} muestreo aleatorio, NearMiss, Tomek Links.\\[2pt]
\good\ Más rápido al entrenar.\\
\bad\ Pierde información. \alert{Inviable en few-shot}.
\end{block}
\end{column}

\begin{column}{0.50\textwidth}
\setbeamercolor{block title}{fg=white,bg=uaiorange}
\begin{block}{Oversampling}
\scriptsize
\textbf{Aumenta} la clase minoritaria.\\[2pt]
\textbf{Ejemplos:} duplicación, \textbf{SMOTE}, ADASYN, EDA, \textbf{LLMs}.\\[2pt]
\good\ Conserva todos los datos.\\
\bad\ Riesgo de overfitting, ruido o divergencia.
\end{block}
\end{column}

\end{columns}

\vspace{0.7em}
\begin{insight}
\scriptsize
SMOTE y los LLMs caen bajo el mismo paraguas conceptual de \textbf{oversampling}, pero operan en planos distintos: \alert{vectores} vs.\ \alert{texto}.\\[2pt]
La pregunta abierta no es cuál elegir, sino \textbf{cómo combinarlos sin perder garantías geométricas}.
\end{insight}

\mensaje{De ahí surge la pregunta central de esta tesis.}
\end{frame}

% --- Slide: Pregunta de investigación ---
\begin{frame}{Pregunta de investigación}
\vfill
\begin{center}
\setbeamercolor{block title}{fg=white,bg=uaibluedark}
\begin{tcolorbox}[colback=white, colframe=uaibluedark, coltitle=white,
                  title=\centering\textbf{Pregunta central},
                  width=0.9\linewidth, halign=center, arc=4pt, boxrule=2pt]
\Large\itshape ¿Puede un sistema de filtrado geométrico\\
post-generación, desacoplado del LLM,\\[0.4em]
\textbf{mejorar significativamente} la aumentación\\
de datos para clasificación de texto\\
con pocos ejemplos por clase?
\end{tcolorbox}
\end{center}
\vfill
\end{frame}


% --- Bloque 4: Estado del Arte ---
% =============================================================================
% Bloque 4: Estado del Arte
% Intro SMOTE/ADASYN + Intro LLMs + Dilema + Timeline + Insights + Posicionamiento
% =============================================================================

\section{Estado del arte}

% --- Slide: Intro SMOTE/ADASYN ---
\begin{frame}{Familia 1 — SMOTE y ADASYN: interpolación geométrica}
\begin{columns}[T]

\begin{column}{0.55\textwidth}
\small
\textbf{SMOTE} \cita{Chawla et al., 2002} y \textbf{ADASYN} \cita{He et al., 2008} son los métodos canónicos de oversampling.

\vspace{0.4em}
\textbf{Idea:}
\begin{itemize}\setlength\itemsep{2pt}
  \item Para cada muestra minoritaria $x_i$, encontrar sus $k$ vecinos más cercanos de la misma clase.
  \item Generar muestras nuevas sobre los segmentos:
\end{itemize}

\vspace{0.2em}
\centering
$x_{\text{new}} = x_i + \lambda \,(x_j - x_i), \quad \lambda \in [0,1]$
\raggedright

\vspace{0.5em}
\textbf{Diferencia clave:}
\begin{itemize}\setlength\itemsep{2pt}
  \item ADASYN \alert{pesa más} las muestras difíciles (rodeadas de otra clase).
\end{itemize}

\vspace{0.4em}
\good\ Geometría controlada. \quad \good\ Garantías formales.\\
\bad\ Genera \alert{vectores}, no texto. \quad \bad\ Solo rellena lo conocido.
\end{column}

\begin{column}{0.45\textwidth}
\centering
\vspace{0.3em}
\begin{tikzpicture}[scale=1.0]
  % Casco convexo
  \fill[methodSMOTE!10] (-1.3,-1.1) -- (1.3,-0.9) -- (1.5,0.7) -- (-0.5,1.3) -- (-1.3,-1.1);

  % Puntos reales
  \foreach \pt in {(-1.3,-1.1),(1.3,-0.9),(1.5,0.7),(-0.5,1.3),(0,0.2)}{\fill[methodSMOTE] \pt circle (3.5pt);}

  % Segmentos de interpolación
  \draw[methodSMOTE!60, dashed] (-1.3,-1.1) -- (1.3,-0.9);
  \draw[methodSMOTE!60, dashed] (-1.3,-1.1) -- (0,0.2);
  \draw[methodSMOTE!60, dashed] (0,0.2) -- (1.5,0.7);
  \draw[methodSMOTE!60, dashed] (-0.5,1.3) -- (0,0.2);

  % Sintéticos (en los segmentos)
  \foreach \pt in {(0,-1),(-0.65,-0.45),(0.75,0.45),(-0.25,0.75)}{\draw[uaiorange, very thick] \pt circle (3pt);}

  \node[font=\tiny, color=uaigray] at (0,-1.8) {{\color{methodSMOTE}\textbullet}\,reales \quad {\color{uaiorange}$\circ$}\,sintéticos};
  \node[font=\tiny\itshape] at (0,-2.1) {Interpolación entre $k$-NN};
\end{tikzpicture}
\end{column}

\end{columns}

\mensaje{Geometría como control: pero la salida es un vector, no una oración.}
\end{frame}

% --- Slide: Familia 2 — Aumentación textual ligera (entre SMOTE y LLMs) ---
\begin{frame}{Familia 2 — Aumentación textual ligera}
\small
Entre la \alert{interpolación geométrica} y la \alert{generación libre} hay un puente histórico: métodos que producen texto con modelos pequeños y operaciones controladas.

\vspace{0.5em}
\begin{columns}[T]

\begin{column}{0.33\textwidth}
\setbeamercolor{block title}{fg=white,bg=uaiblue}
\begin{block}{EDA}
\scriptsize
\cita{Wei \& Zou, 2019}\\[3pt]
4 ops a nivel \textbf{palabra}:\\
sinónimo, inserción, swap, deleción.\\[3pt]
\good\ Muy barato.\\
\bad\ Ruido en oraciones cortas.
\end{block}
\end{column}

\begin{column}{0.33\textwidth}
\setbeamercolor{block title}{fg=white,bg=uaiblue}
\begin{block}{Traducción inversa}
\scriptsize
\cita{Sennrich 2016; Edunov 2018}\\[3pt]
EN $\to$ otro idioma $\to$ EN.\\
Paráfrasis a nivel \textbf{oración}.\\[3pt]
\good\ Cambia sintaxis.\\
\bad\ Pierde matices.
\end{block}
\end{column}

\begin{column}{0.33\textwidth}
\setbeamercolor{block title}{fg=white,bg=uaiblue}
\begin{block}{BERT contextual}
\scriptsize
\cita{Kobayashi 2018; Wu 2019}\\[3pt]
Enmascarar tokens; BERT predice alternativas plausibles.\\[3pt]
\good\ Contexto semántico.\\
\bad\ Solo sustitución 1:1.
\end{block}
\end{column}

\end{columns}

\vspace{0.7em}
\begin{insight}
\scriptsize
Todos producen \textbf{texto}, pero con un \textbf{techo de creatividad bajo}: paráfrasis cercanas al original. Los LLMs rompen ese techo --- y por eso necesitan otro tipo de control.
\end{insight}
\end{frame}

% --- Slide: Familia 3 — LLMs ---
\begin{frame}{Familia 3 — LLMs: generación de texto bajo demanda}
\begin{columns}[T]

\begin{column}{0.58\textwidth}
\small
Pedirle a un LLM (GPT, Gemini, Claude\dots) que genere ejemplos nuevos para cada clase.

\vspace{0.3em}
\textbf{Receta típica} \cita{Whitehouse et al., 2023}:
\begin{enumerate}\setlength\itemsep{1pt}\scriptsize
  \item Prompt con descripción de clase y ejemplos reales.
  \item Pedir $N$ generaciones nuevas.
  \item Añadirlas al entrenamiento.
\end{enumerate}

\vspace{0.3em}
\textbf{Trabajos representativos:}
\begin{itemize}\setlength\itemsep{1pt}\scriptsize
  \item Suhaeni \& Yong (2024): $+12.7\%$ acc.\ en minoritarias.
  \item ImbLLM (2024): control de desbalance vía prompting.
  \item Cegin et al.\ (2025): diversidad guiada por LLM.
\end{itemize}

\vspace{0.4em}
{\scriptsize
\good\ Texto coherente. \quad \good\ Diversidad lingüística.\\
\bad\ \alert{Sin garantías geométricas}. \quad \bad\ Ruido y solapamiento.}
\end{column}

\begin{column}{0.42\textwidth}
\centering
\vspace{0.2em}
\begin{tikzpicture}[every node/.style={font=\scriptsize, align=center}]
  \node[draw=uaibluedark, fill=uaibluelight!15, rounded corners, very thick,
        text width=3.6cm, minimum height=1.0cm] (prompt) at (0, 2.7)
        {\scriptsize\textbf{Prompt}\\[1pt]
         \tiny\itshape ``Genera 50 mensajes de alegría\dots''};
  \draw[->, very thick, uaibluedark] (0, 2.05) -- (0, 1.65);
  \node[draw=methodOurs, fill=methodOurs!15, rounded corners, very thick,
        text width=2.4cm, minimum height=0.7cm] (llm) at (0, 1.2)
        {\textbf{LLM}};
  \draw[->, very thick, methodOurs] (0, 0.8) -- (0, 0.4);
  \node[draw=uaiorange, fill=uaiorange!10, rounded corners, very thick,
        text width=3.6cm, minimum height=1.0cm] (out) at (0, -0.2)
        {\scriptsize\textbf{$N$ oraciones nuevas}\\[1pt]
         \tiny algunas buenas, otras ruido\dots};
\end{tikzpicture}
\end{column}

\end{columns}

\mensaje{Texto sin control: los vectores caen donde quieran.}
\end{frame}

% --- Slide: El dilema (movido desde Introducción) ---
\begin{frame}{El dilema: SMOTE vs.\ LLMs lado a lado}
\begin{columns}[T]

\begin{column}{0.5\textwidth}
\centering
\setbeamercolor{block title}{fg=white,bg=methodSMOTE}
\begin{block}{SMOTE / interpolación}
\small
\good\ Geometría controlada\\
\good\ Garantías formales\\[2pt]
\bad\ Vectores sin texto coherente\\
\bad\ Solo rellena lo conocido
\end{block}
\vspace{0.4em}
\begin{tikzpicture}[scale=0.85]
  \fill[methodSMOTE!20] (0,0) circle (1.4);
  \foreach \pt in {(0,0.4),(0.7,0.2),(-0.5,0.6),(0.3,-0.5),(-0.7,-0.3),(0.6,-0.7)} {\fill[methodSMOTE] \pt circle (2pt);}
  \foreach \pt in {(0.4,0),(-0.2,0.1),(0.1,0.5),(-0.4,-0.5),(0.45,-0.2)} {\draw[uaiorange,thick] \pt circle (2pt);}
  \node[font=\tiny] at (0,-1.7) {Interpolación lineal};
\end{tikzpicture}
\end{column}

\begin{column}{0.5\textwidth}
\centering
\setbeamercolor{block title}{fg=white,bg=methodOurs}
\begin{block}{LLM (Gemini, GPT, Claude\dots)}
\small
\good\ Texto fluido y coherente\\
\good\ Diversidad lingüística\\[2pt]
\bad\ Geometría descontrolada\\
\bad\ Ruido y solapamiento entre clases
\end{block}
\vspace{0.4em}
\begin{tikzpicture}[scale=0.85]
  \fill[methodOurs!15] (0,0) circle (1.6);
  \foreach \pt in {(0,0.4),(0.7,0.2),(-0.5,0.6),(0.3,-0.5)} {\fill[methodOurs] \pt circle (2pt);}
  \foreach \pt in {(0.4,0),(-0.2,0.1),(0.1,0.5),(1.4,0.3),(-1.2,-0.4),(0.9,-1.0),(-0.8,1.1)} {\draw[uaiorange,thick] \pt circle (2pt);}
  \node[font=\tiny] at (0,-1.9) {Generación abierta};
\end{tikzpicture}
\end{column}

\end{columns}

\mensaje{¿Podemos quedarnos con lo mejor de ambos mundos?}
\end{frame}

% --- Slide: Timeline de 4 paradigmas ---
\begin{frame}{Cuatro paradigmas, una década de evolución}
\centering
\begin{tikzpicture}[every node/.style={font=\scriptsize}]
  % Línea de tiempo
  \draw[->, very thick, uaiblue!70] (0,0) -- (12,0) node[right, font=\scriptsize]{año};
  \foreach \x/\year in {0.5/2002, 3.5/2019, 6.5/2023, 9.5/2025}{
    \draw[uaiblue] (\x,-0.1) -- (\x,0.1);
    \node[below, font=\tiny] at (\x,-0.15) {\year};
  }

  % Paradigma 1
  \node[draw=uaiblue, fill=uaibluelight!15, rounded corners, text width=2.4cm, align=center,
        minimum height=1.2cm, font=\tiny] at (1.5, 1.5) (p1)
        {\textbf{Geométrico clásico}\\SMOTE, ADASYN, EDA};
  \draw[->, uaiblue!50] (0.5,0) -- (p1);

  % Paradigma 2
  \node[draw=uaiblue, fill=uaibluelight!15, rounded corners, text width=2.4cm, align=center,
        minimum height=1.2cm, font=\tiny] at (4.5, 1.5) (p2)
        {\textbf{Embeddings profundos}\\mixup, MixText, SMOTExT};
  \draw[->, uaiblue!50] (3.5,0) -- (p2);

  % Paradigma 3
  \node[draw=uaiblue, fill=uaibluelight!15, rounded corners, text width=2.4cm, align=center,
        minimum height=1.2cm, font=\tiny] at (7.5, 1.5) (p3)
        {\textbf{Generación LLM}\\GPT-3/4, Gemini, Claude};
  \draw[->, uaiblue!50] (6.5,0) -- (p3);

  % Paradigma 4
  \node[draw=uaiorange, very thick, fill=uaiorange!15, rounded corners, text width=2.4cm, align=center,
        minimum height=1.2cm, font=\tiny] at (10.5, 1.5) (p4)
        {\textbf{Híbridos + filtros}\\\textit{esta tesis}};
  \draw[->, uaiorange] (9.5,0) -- (p4);

  % Notas debajo
  \node[font=\tiny, text width=2.6cm, align=center] at (1.5,-1) {Chawla 2002\\Wei \& Zou 2019};
  \node[font=\tiny, text width=2.6cm, align=center] at (4.5,-1) {Zhang 2018\\Bayer 2023};
  \node[font=\tiny, text width=2.6cm, align=center] at (7.5,-1) {Suhaeni \& Yong\\\textbf{+12.7\%} acc};
  \node[font=\tiny, text width=2.6cm, align=center] at (10.5,-1) {Cegin 2025, ImbLLM\\Springer 2025};
\end{tikzpicture}

\vfill
\mensaje{Cada paradigma resuelve un problema y crea otro nuevo.}
\end{frame}

% --- Slide: Tres deficiencias (después del Timeline) ---
\begin{frame}{Tres deficiencias en los métodos actuales}
\small
\begin{columns}[T]

\begin{column}{0.33\textwidth}
\setbeamercolor{block title}{fg=white,bg=uaiorange}
\begin{block}{1. Sin validación}
Las muestras del LLM se aceptan \textbf{todas}, sin verificar si caen en la región correcta del espacio.
\vspace{0.3em}\\
{\scriptsize Texto gramatical $\neq$ texto en la zona correcta.}
\end{block}
\end{column}

\begin{column}{0.33\textwidth}
\setbeamercolor{block title}{fg=white,bg=uaiamber}
\begin{block}{2. Peso uniforme}
Todas las muestras pesan igual en el clasificador, ignorando la \textbf{calidad geométrica} de cada una.
\vspace{0.3em}\\
{\scriptsize Una en zona densa = una en zona dispersa.}
\end{block}
\end{column}

\begin{column}{0.33\textwidth}
\setbeamercolor{block title}{fg=white,bg=uaiblue}
\begin{block}{3. Un solo dataset}
Casi todos los estudios validan en \textbf{un único} corpus o tarea.
\vspace{0.3em}\\
{\scriptsize ¿Generaliza el filtrado entre dominios y tareas?}
\end{block}
\end{column}

\end{columns}

\vfill
\begin{tikzpicture}[overlay,remember picture]
\end{tikzpicture}

\mensaje{La señal geométrica está disponible y nadie la está usando.}
\end{frame}

% --- Slide: Gap metodológico + idea conceptual ---
\begin{frame}{El gap metodológico: filtrar la salida del LLM}
\begin{columns}[T]

\begin{column}{0.52\textwidth}
\small
Las deficiencias anteriores apuntan a una intervención conceptualmente simple, no estudiada sistemáticamente.

\vspace{0.4em}
\textbf{La idea:}
\begin{enumerate}\setlength\itemsep{1pt}\scriptsize
  \item Generar $N$ muestras con el LLM (como hoy).
  \item Proyectarlas al espacio de embeddings.
  \item Medir qué tan cerca caen del \textbf{cluster real} de su clase.
  \item \alert{Rechazar las que se alejan demasiado}.
\end{enumerate}

\vspace{0.4em}
\textbf{¿Por qué tiene sentido?}
\begin{itemize}\setlength\itemsep{1pt}\scriptsize
  \item La geometría \alert{ya está ahí}: usamos embeddings para el clasificador.
  \item Cero costo extra de generación; el filtro opera \emph{a posteriori}.
  \item Agnóstico al LLM concreto.
\end{itemize}
\end{column}

\begin{column}{0.48\textwidth}
\centering
\vspace{0.2em}
\begin{tikzpicture}[scale=0.85]
  % Cluster real (fondo)
  \fill[methodSMOTE!18] (0,0) circle (1.5);

  % Frontera del filtro (dashed)
  \draw[uaibluedark, dashed, thick] (0,0) circle (1.9);

  % Puntos reales (sólidos)
  \foreach \pt in {(0,0.2),(0.4,0.4),(-0.5,0.3),(0.6,-0.4),(-0.3,-0.5),(0.2,0.7),(-0.6,-0.2),(0.7,0.1)} {\fill[methodSMOTE] \pt circle (3pt);}

  % LLM aceptados (verdes, dentro del filtro)
  \foreach \pt in {(0.9,0.5),(-0.8,-0.7),(0.3,1.0),(-0.4,1.1),(1.0,-0.5),(-1.1,0.4)} {\draw[methodOurs, very thick] \pt circle (3pt);}

  % LLM rechazados (rojos, fuera del filtro)
  \foreach \pt in {(2.4,0.7),(-2.3,0.9),(2.0,-1.4),(-1.7,-2.0),(0.1,2.3),(2.5,-0.3)} {\draw[uaiorange, very thick] \pt circle (3pt);}

  % Leyenda compacta
  \node[font=\tiny, color=methodSMOTE, anchor=west] at (-2.8,3.0) {{\fontsize{6}{6}\selectfont$\bullet$}\,reales};
  \node[font=\tiny, color=methodOurs, anchor=west] at (-2.8,2.7) {$\circ$\,aceptado};
  \node[font=\tiny, color=uaiorange, anchor=west] at (-2.8,2.4) {$\circ$\,rechazado};
  \node[font=\tiny, color=uaibluedark, anchor=west] at (-2.8,2.1) {- -\,filtro};
\end{tikzpicture}
\end{column}

\end{columns}

\mensaje{La señal geométrica ya existe — falta usarla para decidir qué entra al entrenamiento.}
\end{frame}

% --- Slide: Posicionamiento ---
\begin{frame}{Dónde encaja esta propuesta}
\begin{center}
\begin{tikzpicture}[scale=0.72, every node/.style={font=\scriptsize}, transform shape]
  \fill[uaiblue!25, opacity=0.7] (-1.4,0) circle (1.9);
  \fill[uaigreen!25, opacity=0.7] (1.4,0) circle (1.9);
  \fill[uaiamber!25, opacity=0.7] (0,-1.9) circle (1.9);

  \node[font=\bfseries\small, color=uaiblue, text width=2.4cm, align=center] at (-3.2,1.4) {Validación post-generación};
  \node[font=\bfseries\small, color=uaigreen!70!black, text width=2.4cm, align=center] at (3.2,1.4) {Ponderación adaptativa};
  \node[font=\bfseries\small, color=uaiamber!70!black, text width=2.6cm, align=center] at (0,-3.4) {Validación multi-tarea (NER)};

  \node[font=\tiny, color=uaiblue!70!black] at (-2.2,0) {LOF};
  \node[font=\tiny, color=uaigreen!60!black, text width=1.9cm, align=center] at (2.2,0) {Importance Sampling};
  \node[font=\tiny, color=uaiamber!70!black] at (0,-1.3) {MELM 2022};

  \node[draw=uaiorange, very thick, fill=white, rounded corners,
        font=\bfseries\scriptsize, text width=1.9cm, align=center,
        minimum height=0.8cm] at (0,-0.3) {Esta propuesta};
\end{tikzpicture}
\end{center}

\mensaje{Contribución: integrar las tres dimensiones por primera vez.}
\end{frame}


% --- Bloque 5: Hipótesis y Objetivos ---
% =============================================================================
% Bloque 5: Hipótesis y Objetivos (slide 11)
% =============================================================================

\section{Hipótesis y objetivos}

% --- Slide 11: H1 + sub-hipótesis + objetivos ---
\begin{frame}{Hipótesis y objetivos}
\setbeamercolor{block title}{fg=white,bg=uaiorange}
\begin{block}{H1 — Hipótesis principal}
\small
El filtrado geométrico de muestras LLM en el espacio de embeddings produce
mejoras \textbf{estadísticamente significativas} en macro F1 respecto a SMOTE,
en clasificación con pocos ejemplos por clase.
\end{block}

\vspace{0.4em}
\begin{columns}[T]

\begin{column}{0.5\textwidth}
\setbeamercolor{block title}{fg=white,bg=uaiblue}
\begin{block}{Sub-hipótesis verificables}
\scriptsize
\begin{description}
  \item[H1a] Filtrado por distancia $>$ sin filtrar.
  \item[H1b] Ponderación suave $\geq$ filtrado binario.
  \item[H1c] Generaliza a NER sin modificación.
  \item[H1d] Lineales $>$ ensambles de árboles.
\end{description}
\end{block}
\end{column}

\begin{column}{0.5\textwidth}
\setbeamercolor{block title}{fg=white,bg=uaibluedark}
\begin{block}{Objetivos específicos}
\scriptsize
\begin{enumerate}
  \item Diseñar e implementar un sistema de filtros con ponderación suave.
  \item Comparar contra \textbf{9 métodos} (clásicos + modernos) sobre múltiples datasets.
  \item Demostrar generalización a \textbf{NER}.
\end{enumerate}
\end{block}
\end{column}

\end{columns}

\vspace{0.4em}
{\scriptsize\itshape Criterio de aceptación:}
{\scriptsize $\Delta>0$,\quad $p_{\text{Bonf}}<0.05$,\quad win-rate $>50\%$.}
\end{frame}


% --- Bloque 6: Metodología ---
% =============================================================================
% Bloque 6: Metodología (slides 12–16)
% =============================================================================

\section{Metodología}

% --- Slide 12: Pipeline arquitectónico ---
\begin{frame}{Pipeline: del texto real al clasificador entrenado}
\centering
\begin{tikzpicture}[
    every node/.style={font=\scriptsize},
    box/.style={draw=uaiblue, rounded corners, fill=uaibluelight!15,
                minimum width=2cm, minimum height=0.9cm, text width=2cm, align=center},
    box2/.style={draw=uaiorange, very thick, rounded corners, fill=uaiorange!15,
                 minimum width=2cm, minimum height=0.9cm, text width=2cm, align=center},
    box3/.style={draw=uaigreen, very thick, rounded corners, fill=uaigreen!15,
                 minimum width=2cm, minimum height=0.9cm, text width=2cm, align=center},
    arrow/.style={->, >=Stealth, thick, uaiblue!70}
  ]
  % Fila superior
  \node[box]  (real)  at (0,2)    {\textbf{Datos\\reales}\\\tiny $n$-shot 10/25/50};
  \node[box]  (enc)   at (2.6,2)  {\textbf{Encoder}\\\tiny mpnet-base-v2\\\tiny 768 dims};
  \node[box]  (llm)   at (5.2,2)  {\textbf{LLM}\\\tiny Gemini 3 Flash\\\tiny $3\times$ excedente};
  \node[box]  (cands) at (7.8,2)  {\textbf{Candidatas}\\\tiny embeddings};
  \node[box2] (filt)  at (10.4,2) {\textbf{Filtros}\\\tiny geométricos};

  % Fila inferior
  \node[box2] (weight) at (10.4,0)  {\textbf{Pesos}\\\tiny softw.\ $T=0.5$};
  \node[box3] (clf)    at (5.2,0)   {\textbf{Clasificador}\\\tiny LR / SVC / Ridge};
  \node[box]  (eval)   at (0,0)     {\textbf{Macro F1}\\\tiny test fijo};

  % Flechas
  \draw[arrow] (real)  -- (enc);
  \draw[arrow] (enc)   -- (llm);
  \draw[arrow] (llm)   -- (cands);
  \draw[arrow] (cands) -- (filt);
  \draw[arrow] (filt)  -- (weight);
  \draw[arrow] (weight) -- (clf);
  \draw[arrow] (real.south)  .. controls +(0,-0.6) and +(-0.5,0) .. (clf.west);
  \draw[arrow] (clf)   -- (eval);
\end{tikzpicture}

\vspace{0.5em}
\mensaje{Filtrado y ponderación: dos contribuciones desacopladas del LLM.}
\end{frame}

% --- Slide 13: Protocolo experimental (vuelve, antes de los datasets) ---
\begin{frame}{Protocolo experimental: 8.000+ configuraciones}
\small
\begin{columns}[T]

\begin{column}{0.55\textwidth}
\textbf{Datos y modelos:}
\begin{itemize}\setlength\itemsep{0pt}
  \item \textbf{13 datasets} en 9 dominios.
  \item \textbf{$n$-shot:} 10 / 25 / 50 ejemplos por clase.
  \item \textbf{5 clasificadores:} LR, SVC, Ridge, RF, MLP.
  \item \textbf{10 métodos} de aumentación.
  \item \textbf{5 LLMs} (Gemini, GPT-5, Claude, Kimi, GLM).
  \item \textbf{4 modelos} de embeddings.
  \item \textbf{5 semillas:} 42, 123, 456, 789, 1011.
\end{itemize}
\end{column}

\begin{column}{0.45\textwidth}
\textbf{Protocolo estadístico:}
\begin{itemize}\setlength\itemsep{0pt}
  \item Métrica: \textbf{macro F1}.
  \item $\Delta$F1 vs.\ SMOTE (puntos porcentuales).
  \item \textbf{$t$-pareada} con corrección de Bonferroni.
  \item Bootstrap CI 95\% (10\,000 remuestreos).
  \item $d$ de Cohen como tamaño de efecto.
  \item Win-rate como medida de consistencia.
\end{itemize}

\vspace{0.3em}
\setbeamercolor{block title}{fg=white,bg=uaiorange}
\begin{block}{Volumen}
\centering\bfseries\large 8.000+ configuraciones
\end{block}
\end{column}

\end{columns}
\end{frame}

% --- Slide 14: Conjuntos de datos (tabla) ---
\begin{frame}{Conjuntos de datos: 13 corpus en 9 dominios}
\centering
\scriptsize
\renewcommand{\arraystretch}{0.92}
\begin{tabular}{@{}llccr@{}}
\toprule
\textbf{Conjunto} & \textbf{Dominio} & \textbf{Clases} & \textbf{$n$-shots} & \textbf{Test} \\
\midrule
\rowcolor{uaibluelight!12}\multicolumn{5}{l}{\textit{\color{uaibluedark}Clasificación de texto --- experimento principal}} \\
\dataset{sms\_spam}              & Detección de spam        & 2  & 10, 25, 50 & 1.394 \\
\dataset{hate\_speech\_davidson} & Discurso de odio         & 3  & 10, 25, 50 & 4.948 \\
\dataset{20newsgroups}           & Noticias (4 cat.)        & 4  & 10, 25, 50 & 7.532 \\
\dataset{ag\_news}               & Noticias (4 cat.)        & 4  & 10, 25, 50 & 7.600 \\
\dataset{emotion}                & Emociones                & 6  & 10, 25, 50 & 2.000 \\
\dataset{dbpedia14}              & Clasificación de temas   & 14 & 10, 25, 50 & 5.000 \\
\dataset{20newsgroups\_20class}  & Noticias (20 cat.)       & 20 & 10, 25, 50 & 7.532 \\
\midrule
\rowcolor{uaiorange!12}\multicolumn{5}{l}{\textit{\color{uaibluedark}Adicionales --- escalabilidad por número de clases}} \\
\dataset{trec6}                  & Preguntas                & 6   & 10, 25, 50 & 500 \\
\dataset{banking77}              & Intenciones bancarias    & 77  & 10, 25, 50 & 3.080 \\
\dataset{clinc150}               & Intenciones de diálogo   & 150 & 10, 25, 50 & 4.500 \\
\midrule
\rowcolor{uaigreen!12}\multicolumn{5}{l}{\textit{\color{uaibluedark}Reconocimiento de entidades nombradas (NER)}} \\
\dataset{MultiNERD}              & Entidades generales      & 5 tipos & 10, 25, 50 & 10.000 \\
\dataset{WikiANN}                & Entidades Wikipedia      & 3 tipos & 10, 25, 50 & 10.000 \\
\dataset{Few-NERD}               & Entidades detalladas     & 8 tipos & 10, 25, 50 & 10.000 \\
\bottomrule
\end{tabular}

\vspace{0.2em}
{\scriptsize\itshape $n$-shot estricto: 10 / 25 / 50 ejemplos por clase. Rango de complejidad: de binario (\dataset{sms\_spam}) a 150 clases (\dataset{clinc150}).}
\end{frame}

% --- Slide 15: Por qué estos datasets de clasificación ---
\begin{frame}{¿Por qué estos datasets de clasificación?}
\begin{columns}[T]

\begin{column}{0.50\textwidth}
\setbeamercolor{block title}{fg=white,bg=uaibluedark}
\begin{block}{\small Experimento principal --- 7 datasets}
\scriptsize
Buscan \textbf{cobertura en tres ejes}:
\begin{itemize}\setlength\itemsep{0pt}
  \item \textbf{Granularidad:} 2 $\to$ 20 clases.
  \item \textbf{Dominio:} spam, odio, noticias, emoción, ontología.
  \item \textbf{Estilo:} corto/largo, formal/informal, social/curado.
\end{itemize}
\vspace{0.4em}
\begin{itemize}\setlength\itemsep{0pt}
  \item \dataset{sms\_spam} (2): NLP canónico, frase corta.
  \item \dataset{hate\_speech\_davidson} (3): dominio sensible.
  \item \dataset{20newsgroups}, \dataset{ag\_news} (4): noticias.
  \item \dataset{emotion} (6): el ejemplo recurrente.
  \item \dataset{dbpedia14} (14): ontología Wikipedia.
  \item \dataset{20newsgroups\_20class} (20): máxima granularidad.
\end{itemize}
\end{block}
\end{column}

\begin{column}{0.50\textwidth}
\setbeamercolor{block title}{fg=white,bg=uaiorange}
\begin{block}{\small Escalabilidad --- 3 datasets adicionales}
\scriptsize
Responden la pregunta:\\[2pt]
\alert{¿el filtrado sigue funcionando con muchas clases?}
\vspace{0.4em}
\begin{itemize}\setlength\itemsep{0pt}
  \item \dataset{trec6} (6): preguntas TREC, mismo régimen de clases que el principal pero distinto dominio. Sirve de \textbf{control}.
  \item \dataset{banking77} (77): intenciones bancarias, granularidad fina con clases semánticamente cercanas.
  \item \dataset{clinc150} (150): intenciones de diálogo de asistentes virtuales. \textbf{Caso extremo}.
\end{itemize}
\vspace{0.4em}
Permiten medir \textbf{degradación del filtrado} al aumentar el número de clases.
\end{block}
\end{column}

\end{columns}

\mensaje{Diez datasets, dos preguntas: ¿qué tan bueno es el filtrado? ¿hasta dónde escala?}
\end{frame}

% --- Slide 16: ¿Qué es NER y por qué la usamos? ---
\begin{frame}{¿Qué es NER y por qué la usamos?}

\setbeamercolor{block title}{fg=white,bg=uaibluedark}
\begin{block}{\small NER (\textit{Named Entity Recognition})}
\scriptsize
Tarea de NLP que \textbf{detecta nombres propios dentro del texto y los clasifica por tipo}: personas, lugares, organizaciones, fechas, etc.

\vspace{0.3em}
\textit{Ejemplo:} «Llamé a \colorbox{uaiblue!25}{\strut\,María\,} en \colorbox{uaiorange!25}{\strut\,Madrid\,} el \colorbox{uaigreen!25}{\strut\,2 de marzo\,}.»\hspace{0.5em}{\tiny \textcolor{uaiblue}{$\uparrow$PER}\hspace{1em}\textcolor{uaiorange}{$\uparrow$LOC}\hspace{1em}\textcolor{uaigreen}{$\uparrow$DATE}}
\end{block}

\vspace{0.4em}
\begin{columns}[T]

\begin{column}{0.55\textwidth}
\small
\textbf{¿Por qué la usamos como prueba?}
\begin{itemize}\setlength\itemsep{1pt}\scriptsize
  \item Construimos el filtro para \textbf{clasificación}.
  \item NER es una tarea NLP \textbf{distinta}: en vez de etiquetar la oración entera, se etiquetan palabras dentro.
  \item Si el mismo filtro funciona en NER \textbf{sin cambiarlo} $\Rightarrow$ el filtro \alert{no depende de la tarea}.
\end{itemize}

\vspace{0.2em}
{\scriptsize Es la prueba más fuerte de que el método \textbf{generaliza}.}
\end{column}

\begin{column}{0.45\textwidth}
\setbeamercolor{block title}{fg=white,bg=uaigreen}
\begin{block}{\small 3 corpus, distinta complejidad}
\scriptsize
\begin{itemize}\setlength\itemsep{1pt}
  \item \dataset{WikiANN} --- 3 tipos (PER, ORG, LOC). \textbf{Fácil}.
  \item \dataset{MultiNERD} --- 5 tipos (+ DATE, MISC). Medio.
  \item \dataset{Few-NERD} --- 8 tipos jerárquicos. \textbf{Difícil}.
\end{itemize}
\end{block}
\end{column}

\end{columns}

\mensaje{Mismo filtro, otra tarea: si funciona acá, sirve en general.}
\end{frame}

% --- Slide 17: Los 5 filtros ---
\begin{frame}{Cinco estrategias de filtrado evaluadas}
\small
\begin{tabular}{@{}p{2.6cm}p{4.6cm}p{4.6cm}@{}}
\toprule
\textbf{Filtro} & \textbf{Criterio} & \textbf{Parámetros} \\
\midrule
\textit{None} (control)        & Sin filtrado, todas las muestras LLM. & --- \\[2pt]
\textit{LOF}                   & Densidad local (Local Outlier Factor). & $k=20$, threshold $=0$ \\[2pt]
\textit{Cascade L1} \alert{$\star$} & Distancia euclidiana al ancla. Top-$N$ por ranking. & $k=10$, niveles 1--4 \\[2pt]
\textit{Combined}              & LOF \textbf{Y} similitud coseno. & $\tau_{\cos}=0.5$ \\[2pt]
\textit{Embedding-Guided}      & Cobertura $0.6$ + calidad $0.4$, voraz. & multi-criterio \\
\bottomrule
\end{tabular}

\vspace{0.6em}
{\scriptsize\itshape\alert{$\star$} Cascade nivel 1 = mejor configuración encontrada.}

\vfill
\mensaje{Cinco diseños, una pregunta: ¿filtros simples o complejos?}
\end{frame}

% --- Slide 14: Cascade — pieza central ---
\begin{frame}{Cascada de filtros: niveles 1--4}
\begin{columns}[T]

\begin{column}{0.55\textwidth}
\small
\textbf{Score compuesto} (media geométrica):
\[
s(x) = \left(\prod_{i=1}^{L} s_i(x)\right)^{1/L}
\]
\begin{description}\scriptsize
  \item[Nivel 1 \alert{$\star$}] Distancia euclidiana al ancla.
  \item[Nivel 2] Similitud coseno.
  \item[Nivel 3] Pureza KNN ($k$=10).
  \item[Nivel 4] Distancia al centroide.
\end{description}

\vspace{0.3em}
\begin{insight}
\scriptsize Con vectores L2-normalizados, niveles 1 y 2 son redundantes; los niveles 3--4 introducen ruido en pocos ejemplos.
\end{insight}
\end{column}

\begin{column}{0.45\textwidth}
\centering
\begin{tikzpicture}[every node/.style={font=\tiny}]
  % Pirámide de niveles
  \fill[uaiorange!40] (-2.5,2) rectangle (2.5,2.7);
  \node[font=\scriptsize\bfseries] at (0,2.35) {Nivel 1: Distancia $\star$};

  \fill[uaibluelight!30] (-2.2,1.2) rectangle (2.2,1.9);
  \node[font=\scriptsize] at (0,1.55) {Nivel 2: Similitud};

  \fill[uaibluelight!20] (-1.9,0.4) rectangle (1.9,1.1);
  \node[font=\scriptsize] at (0,0.75) {Nivel 3: KNN-purity};

  \fill[uaibluelight!10] (-1.6,-0.4) rectangle (1.6,0.3);
  \node[font=\scriptsize] at (0,-0.05) {Nivel 4: Centroide};

  % Flecha de score
  \draw[->, very thick, uaiorange] (-3.2,-0.4) -- (-3.2,2.7);
  \node[font=\tiny, rotate=90] at (-3.6,1.15) {score $\uparrow$};
\end{tikzpicture}\\[4pt]
{\scriptsize Selecciona top-$N$ por ranking,\\ \textbf{sin umbral binario}.}
\end{column}

\end{columns}
\end{frame}

% --- Slide 15: Ponderación suave ---
\begin{frame}{Ponderación suave: la 2.\textsuperscript{a} contribución}
\begin{columns}[T]

\begin{column}{0.5\textwidth}
\small
En vez de \textbf{aceptar / rechazar}, asignamos un \textit{peso} a cada muestra.

\vspace{0.4em}
\textbf{Eq. 1 — Normalización min-max:}
\[
\tilde{s}(x) = \frac{s(x) - s_{\min}}{s_{\max} - s_{\min}}
\]
\textbf{Eq. 2 — Escalado por temperatura:}
\[
w(x) = \max\!\left(w_{\min},\ \tilde{s}(x)^{\,1/T}\right)
\]

{\scriptsize $T=0.5$, $w_{\min}=0$, datos reales $\Rightarrow$ peso 1.}
\end{column}

\begin{column}{0.5\textwidth}
\centering
\begin{tikzpicture}[every node/.style={font=\tiny}]
  % Eje
  \draw[->,uaigray] (0,0) -- (4.5,0) node[right]{$\tilde{s}(x)$};
  \draw[->,uaigray] (0,0) -- (0,3.2) node[above]{$w(x)$};
  \node[font=\tiny] at (4,-0.25) {1};
  \node[font=\tiny] at (-0.2,3) {1};

  % Curva T=0.5 (peaked) — w = s^2
  \draw[very thick, uaiorange, domain=0:4, samples=40, smooth] plot (\x, {3*(\x/4)^2});
  \node[uaiorange, font=\scriptsize] at (3.6,2.6) {$T=0.5$};

  % Curva T=1 (linear)
  \draw[thick, uaiblue, dashed, domain=0:4, samples=40] plot (\x, {3*(\x/4)});
  \node[uaiblue, font=\scriptsize] at (3.6,1.5) {$T=1$};

  \node[align=center, font=\tiny] at (2.2,-0.7) {Score normalizado};
\end{tikzpicture}\\[4pt]
{\scriptsize $T<1$ amplifica diferencias entre buenas y malas.}
\end{column}

\end{columns}

\mensaje{Las muestras se entregan al clasificador como \texttt{sample\_weight}.}
\end{frame}



% --- Bloque 7: Resultados ---
% =============================================================================
% Bloque 7: Resultados (slides 17–24) — corazón de la defensa
% =============================================================================

\section{Resultados}

% --- Slide 17: HEADLINE ---
\begin{frame}{Resultado principal}
\begin{columns}[T]

\begin{column}{0.42\textwidth}
\vspace{0.4em}
\begin{headline}
\cifra{$+2.25$ pp}{vs.\ SMOTE\\ (macro F1)}
\end{headline}

\vspace{0.5em}
\small
\begin{itemize}\setlength\itemsep{2pt}
  \item $p < 0{.}0001$ \cita{Bonferroni}
  \item $d$ de Cohen $= 0{.}74$ \cita{efecto medio-grande}
  \item Win-rate $= 83{.}8\%$
  \item $N = 3{.}675$ configuraciones
  \item Clasificación absoluta: \textbf{73.49\%} vs.\ 71.24\%
\end{itemize}
\end{column}

\begin{column}{0.58\textwidth}
\centering
\includegraphics[width=\linewidth, height=6.5cm, keepaspectratio]{fig_boxplot_deltas.png}\\[2pt]
{\scriptsize Distribución de $\Delta$F1 por método.}
\end{column}

\end{columns}

\mensaje{El método propuesto supera todas las líneas base, incluidas las modernas.}
\end{frame}

% --- Slide 18a: Antes vs después del filtrado ---
\begin{frame}{Visualización t-SNE: el filtro en acción --- antes vs después}
\centering
\includegraphics[width=0.96\linewidth, height=5.4cm, keepaspectratio]{fig_tsne_before_after_20newsgroups.png}

\vspace{0.1em}
{\scriptsize \dataset{20newsgroups} (4 clases, 10-shot). \textbf{Izq.:} candidatas del LLM (triángulos tenues). \textbf{Der.:} solo las que pasaron el filtro. Círculos = reales, estrellas = centroides.}

\mensaje{El filtro elimina las candidatas en zonas de cruce o periféricas; conserva las que refuerzan el cluster.}
\end{frame}

% --- Slide 18b: t-SNE filtrado vs SMOTE ---
\begin{frame}{Visualización t-SNE: SMOTE vs filtro}
\centering
\includegraphics[width=0.85\linewidth, height=5.8cm, keepaspectratio]{fig_tsne_filtered_vs_smote_emotion.png}

\vspace{0.2em}
{\scriptsize Dataset \dataset{emotion} (6 clases, 10-shot): triángulos = LLM filtrado, cuadrados = SMOTE.}

\mensaje{El filtro rechaza candidatas en zonas de cruce y aporta diversidad real.}
\end{frame}

% --- Slide 19: F1 vs n-shot ---
\begin{frame}{Retornos decrecientes con más datos}
\begin{columns}[T]

\begin{column}{0.55\textwidth}
\centering
\includegraphics[width=\linewidth, height=5.8cm, keepaspectratio]{fig_delta_by_nshot.png}\\[2pt]
{\scriptsize $\Delta$F1 por valor de $n$-shot.}
\end{column}

\begin{column}{0.42\textwidth}
\small
\textbf{El método brilla donde más se necesita:}
\vspace{0.3em}

\begin{tabular}{@{}lcc@{}}
\toprule
$n$-shot & $\Delta$F1 & $d$ \\
\midrule
\textbf{10} & \alert{$+4{.}89$} & 1.48 \\
25          & $+2{.}17$         & 1.57 \\
50          & $+0{.}76$         & 0.80 \\
\bottomrule
\end{tabular}

\vspace{0.5em}
\begin{insight}
\scriptsize Cuanto menos datos hay, más útil es el filtrado.
\end{insight}
\end{column}

\end{columns}

\mensaje{Caso de uso ideal: 10 ejemplos por clase.}
\end{frame}

% --- Slide 20: Por clasificador ---
\begin{frame}{El clasificador importa: lineales ganan}
\begin{columns}[T]

\begin{column}{0.55\textwidth}
\centering
\includegraphics[width=\linewidth, height=5.5cm, keepaspectratio]{fig_classifier_comparison.png}
\end{column}

\begin{column}{0.42\textwidth}
\small
\begin{tabular}{@{}lr@{\hspace{6pt}}r@{}}
\toprule
Clasificador & $\Delta$F1 & $d$ \\
\midrule
SVC (lineal)        & \alert{$+2{.}99$} & 1.02 \\
Ridge               & $+2{.}69$         & 1.04 \\
MLP                 & $+2{.}16$         & 0.73 \\
Logistic Regression & $+1{.}98$         & 0.82 \\
Random Forest       & $+1{.}42$         & n.s. \\
\bottomrule
\end{tabular}

\vspace{0.5em}
{\scriptsize H1d confirmada: \textbf{ensambles de árboles}\\ no aprovechan los pesos suaves.}
\end{column}

\end{columns}

\mensaje{Modelos lineales explotan mejor la señal geométrica.}
\end{frame}

% --- Slide 21: Sweet spot por número de clases ---
\begin{frame}{¿Cuándo conviene aplicarlo? Por número de clases}
\begin{columns}[T]

\begin{column}{0.5\textwidth}
\centering
\includegraphics[width=\linewidth, height=5.5cm, keepaspectratio]{fig_heatmap_methods.png}
\end{column}

\begin{column}{0.5\textwidth}
\small
\textbf{Sweet spot: pocas clases.}
\vspace{0.2em}

\begin{tabular}{@{}lr@{}}
\toprule
\# Clases (dataset)     & $\Delta$F1 \\
\midrule
6 (\dataset{trec6, emotion}) & \alert{$+5{.}42$} \\
4 (\dataset{20news, ag})     & $+2{.}80$ \\
2 (binarios)                 & $+3{.}18$ \\
77 (\dataset{banking77})     & $-0{.}99$ \\
150 (\dataset{clinc150})     & $+0{.}01$ \\
\bottomrule
\end{tabular}

\vspace{0.4em}
{\scriptsize Spearman $\rho = -0{.}260$, $p=0{.}003$.}
\end{column}

\end{columns}

\mensaje{Saturación con $\geq 77$ clases: el espacio se fragmenta.}
\end{frame}

% --- Slide 22: Hallazgo contra-intuitivo ---
\begin{frame}{Hallazgo contra-intuitivo: simple $>$ complejo}
\begin{columns}[T]

\begin{column}{0.5\textwidth}
\centering
\begin{tikzpicture}[every node/.style={font=\scriptsize}, scale=0.7]
  % Línea vertical en 0 (referencia SMOTE)
  \draw[uaigray, dashed, thick] (0,-0.3) -- (0,4.8);
  \node[below, font=\tiny] at (0,-0.35) {SMOTE = 0};

  % Eje horizontal abajo
  \draw[->,thick,uaigray] (-3.7,-0.3) -- (2.2,-0.3) node[right, font=\tiny]{$\Delta$F1 (pp)};
  \foreach \x in {-3,-2,-1,1,2}{
    \draw[uaigray!50] (\x,-0.22) -- (\x,-0.38);
  }

  % Bars (positive)
  \fill[methodOurs] (0,4.1) rectangle (1.45,4.5);
  \node[right, font=\tiny] at (1.55,4.3) {\textbf{Cascade L1 +1.45}};

  \fill[methodOurs!70] (0,3.4) rectangle (0.88,3.8);
  \node[right, font=\tiny] at (0.98,3.6) {LOF +0.88};

  \fill[methodOther] (0,2.7) rectangle (0.72,3.1);
  \node[right, font=\tiny] at (0.82,2.9) {None +0.72};

  % Bars (negative)
  \fill[methodWorse] (-1.29,2.0) rectangle (0,2.4);
  \node[left, font=\tiny] at (-1.39,2.2) {Emb.-Guided $-1.29$};

  \fill[methodWorse] (-3.28,1.3) rectangle (0,1.7);
  \node[left, font=\tiny, color=uaiorange] at (-3.38,1.5) {\textbf{Combined $-3.28$}};
\end{tikzpicture}
\end{column}

\begin{column}{0.5\textwidth}
\setbeamercolor{block title}{fg=white,bg=uaiamber}
\begin{block}{Insight clave}
Filtros multi-criterio \textbf{degradan} el rendimiento.
\vspace{0.4em}
\begin{itemize}\setlength\itemsep{0pt}\scriptsize
  \item Distancia y similitud son redundantes en L2-norm.
  \item LOF inestable con $n=10$.
  \item Restricciones conjuntivas matan la diversidad.
\end{itemize}
\end{block}

\vspace{0.4em}
\begin{insight}
\scriptsize\centering
Más filtrado $\neq$ mejor filtrado.\\Generar muchos, filtrar poco.
\end{insight}
\end{column}

\end{columns}
\end{frame}

% --- Slide 23: Robustez multi-LLM y multi-embedding ---
\begin{frame}{Robustez: ¿depende del LLM o del embedding?}
\begin{columns}[T]

\begin{column}{0.5\textwidth}
\centering
\includegraphics[width=\linewidth, height=4.6cm, keepaspectratio]{fig_multi_llm_comparison.png}\\[2pt]
\textbf{5 LLMs probados}\\
{\scriptsize Gemini, GPT-5, Claude, Kimi, GLM}\\
{\scriptsize $\Delta$F1 entre $+1{.}49$ y $+2{.}31$ pp}\\
{\scriptsize Friedman $p=0{.}086$ \cita{no significativo}}
\end{column}

\begin{column}{0.5\textwidth}
\centering
\setbeamercolor{block title}{fg=white,bg=uaiblue}
\begin{block}{4 modelos de embeddings}
\scriptsize
\begin{tabular}{@{}lr@{}}
\toprule
Modelo & $\Delta$F1 \\
\midrule
mpnet-base-v2 (768d)     & $+2{.}74$ \\
bge-large (1024d)        & $+2{.}04$ \\
bge-small (384d)         & $+2{.}00$ \\
e5-large (1024d)         & $+0{.}97$ \\
\bottomrule
\end{tabular}
\vspace{0.3em}\\
{\scriptsize 90.5\% acuerdo entre modelos.}
\end{block}
\end{column}

\end{columns}

\mensaje{El método es agnóstico al LLM y al modelo de embeddings.}
\end{frame}

% --- Slide 24: NER ---
\begin{frame}{Generalización a NER: \textit{el filtro funciona en otra tarea}}
\begin{columns}[T]

\begin{column}{0.45\textwidth}
\small
\textbf{Reconocimiento de entidades nombradas:}
\begin{itemize}\setlength\itemsep{0pt}
  \item 3 corpus: \dataset{MultiNERD}, \dataset{WikiANN}, \dataset{Few-NERD}
  \item Adaptación: clase = tipo de entidad dominante por oración.
  \item \textbf{Filtros sin modificación.}
\end{itemize}

\vspace{0.4em}
\begin{headline}
\cifra{$+9{.}61$ pp}{Cascade L1\\ ($p=0.0003$, $d=2.05$)}
\end{headline}
\end{column}

\begin{column}{0.55\textwidth}
\centering
\begin{tikzpicture}[every node/.style={font=\scriptsize}, x=0.32cm, y=0.5cm]
  % Eje horizontal y vertical
  \draw[->,thick,uaigray] (0,0) -- (16.5,0) node[right, font=\tiny]{$\Delta$F1 (pp)};
  \draw[uaigray] (0,0) -- (0,7.5);
  \foreach \x in {0,3,6,9,12,15}{
    \draw[uaigray!50] (\x,0.15) -- (\x,-0.15) node[below, font=\tiny]{\x};
  }

  % Cascade L1
  \fill[methodOurs] (0,6) rectangle (9.61,7);
  \node[right, font=\scriptsize\bfseries] at (9.81,6.5) {Cascade L1 +9.61};

  % LOF
  \fill[methodOurs!80] (0,4.5) rectangle (9.28,5.5);
  \node[right, font=\scriptsize] at (9.48,5.0) {LOF +9.28};

  % None: solo línea de referencia (sin barra)
  \draw[methodSMOTE, very thick] (-0.3,3) -- (0.3,4);
  \node[right, font=\scriptsize, color=methodSMOTE] at (0.5,3.5) {None (referencia $= 0$)};

  % Combined: marca de rechazo (no medido)
  \node[right, font=\scriptsize, color=uaiorange] at (0.2,2) {\textbf{$\times$} Combined: 0\% aceptación};

  % Anotación 10-shot
  \node[font=\tiny, color=uaigreen!50!black] at (8,0.8) {(10-shot: $+15.15$ pp)};
\end{tikzpicture}
\end{column}

\end{columns}

\mensaje{La calidad geométrica es agnóstica a la tarea.}
\end{frame}


% --- Bloque 8: Conclusiones y cierre ---
% =============================================================================
% Bloque 8: Discusión, Utilidad Real, Cierre (slides 25–27)
% =============================================================================

\section{Discusión y aplicaciones}

% --- Slide 25: Contribuciones agrupadas ---
\begin{frame}{Contribuciones de la tesis}
\small
\begin{columns}[T]

\begin{column}{0.33\textwidth}
\setbeamercolor{block title}{fg=white,bg=uaiblue}
\begin{block}{Metodológica}
\scriptsize
\begin{itemize}\setlength\itemsep{0pt}
  \item Sistema de filtros geométricos multi-nivel.
  \item Ponderación suave por temperatura.
  \item Desacoplado del LLM.
\end{itemize}
\end{block}
\end{column}

\begin{column}{0.33\textwidth}
\setbeamercolor{block title}{fg=white,bg=uaiorange}
\begin{block}{Empírica}
\scriptsize
\begin{itemize}\setlength\itemsep{0pt}
  \item +2.25 pp vs.\ SMOTE\\(\textit{p}\,$<$\,0.0001).
  \item Supera a 9 baselines\\incluidos T5 y BERT.
  \item Robusto a 5 LLMs y 4 embeddings.
\end{itemize}
\end{block}
\end{column}

\begin{column}{0.33\textwidth}
\setbeamercolor{block title}{fg=white,bg=uaigreen}
\begin{block}{Generalización}
\scriptsize
\begin{itemize}\setlength\itemsep{0pt}
  \item NER: +9.61 pp sin modificar nada.
  \item Hallazgo: simple $>$ complejo.
  \item Beneficio en clases difíciles: +10.44 pp.
\end{itemize}
\end{block}
\end{column}

\end{columns}

\vspace{0.5em}
\begin{insight}
\scriptsize\centering
\textbf{Insight central:} la información geométrica del espacio de embeddings\\
captura calidad de manera agnóstica a la tarea, siempre que el filtro sea simple.
\end{insight}
\end{frame}

% --- Slide 26: Utilidad real / aplicaciones ---
\begin{frame}{¿Cuándo y cómo usarlo? Utilidad práctica}
\small
\begin{columns}[T]

\begin{column}{0.55\textwidth}
\textbf{Receta de aplicación:}
\begin{enumerate}\setlength\itemsep{1pt}
  \item Tienes \alert{10--25 ejemplos por clase}.
  \item Tienes \alert{4--10 clases}.
  \item Generas $3\times$ candidatas con un LLM.
  \item Filtras con cascade nivel 1 (distancia).
  \item Ponderas suavemente y entrenas un\\
        clasificador \textbf{lineal}.
\end{enumerate}

\vspace{0.4em}
\setbeamercolor{block title}{fg=white,bg=uaiorange}
\begin{block}{Costo real}
$\sim 10$--$50$ llamadas LLM por clase.\\
A precios actuales: \textbf{centavos por dataset}.
\end{block}
\end{column}

\begin{column}{0.45\textwidth}
\textbf{Casos de aplicación inmediata:}
\begin{itemize}\setlength\itemsep{1pt}
  \item Detectar contenido tóxico nuevo.
  \item Triaje clínico con notas escasas.
  \item Categorización de tickets de soporte.
  \item Monitoreo de marca / tópicos emergentes.
  \item Clasificación de eventos en logs.
\end{itemize}

\vspace{0.4em}
\textbf{Quién no debería usarlo:}
\begin{itemize}\setlength\itemsep{1pt}
  \item Más de 50 ejemplos por clase: SMOTE basta.
  \item $\geq 77$ clases: el espacio se fragmenta.
  \item Modelos de árbol como única opción.
\end{itemize}
\end{column}

\end{columns}
\end{frame}

% --- Slide 27: Limitaciones, futuro y cierre ---
\begin{frame}{Limitaciones, trabajo futuro y cierre}
\small
\begin{columns}[T]

\begin{column}{0.5\textwidth}
\setbeamercolor{block title}{fg=white,bg=uaiamber}
\begin{block}{Limitaciones honestas}
\scriptsize
\begin{itemize}\setlength\itemsep{0pt}
  \item Evaluación solo en \textbf{inglés}.
  \item Caching reduce varianza\\\cita{$+1{.}79$pp persistente en stochastic clf}.
  \item Pierde efectividad con $\geq 77$ clases.
  \item LLMs comerciales filtran contenido sensible (hate-speech).
\end{itemize}
\end{block}
\end{column}

\begin{column}{0.5\textwidth}
\setbeamercolor{block title}{fg=white,bg=uaiblue}
\begin{block}{Trabajo futuro}
\scriptsize
\begin{itemize}\setlength\itemsep{0pt}
  \item Embeddings y LLMs \textbf{multilingües}.
  \item LLMs locales (Llama, Mistral) sin filtros de contenido.
  \item Temperatura $T$ \textbf{adaptativa por clase}.
  \item Presupuesto generativo adaptativo.
  \item Otras tareas: QA, resumen, generación de preguntas.
\end{itemize}
\end{block}
\end{column}

\end{columns}

\vspace{0.6em}
\begin{center}
\setbeamercolor{block title}{fg=white,bg=uaibluedark}
\begin{tcolorbox}[colback=uaibluedark, colframe=uaibluedark, coltext=white,
                  width=0.85\linewidth, halign=center, arc=4pt]
\textbf{Gracias.}\\[0.2em]
{\small Preguntas, comentarios, contraejemplos: bienvenidos.}
\end{tcolorbox}
\end{center}
\end{frame}


% --- Bloque 9: Backup (slides para Q&A) ---
\appendix
% =============================================================================
% Bloque 9: Backup (slides B1–B5) — para Q&A
% =============================================================================

% --- B1: Validación de SMOTE ---
\begin{frame}{Backup B1: SMOTE — validación de implementación}
\small
\begin{columns}[T]

\begin{column}{0.55\textwidth}
\textbf{Crítica anticipada del comité:}
\begin{quote}\itshape\scriptsize
``¿Su SMOTE \textit{dummy-binary} es equivalente al SMOTE estándar binary-mask?''
\end{quote}

\textbf{Validación:}
\begin{itemize}\setlength\itemsep{0pt}
  \item 27 comparaciones lado a lado.
  \item Diferencia máxima: $0{.}0000$ pp.
  \item F1 idénticos en todos los datasets.
\end{itemize}

\vspace{0.3em}
\begin{insight}
\scriptsize Línea base SMOTE \textbf{robusta y reproducible}.
\end{insight}
\end{column}

\begin{column}{0.45\textwidth}
\centering
\setbeamercolor{block title}{fg=white,bg=uaigreen}
\begin{block}{Resumen}
\scriptsize
\begin{tabular}{@{}lr@{}}
\toprule
Implementación   & F1 medio \\
\midrule
dummy-binary     & 0.7124 \\
binary-mask std. & 0.7124 \\
\midrule
Diferencia       & \textbf{0.0000} \\
\bottomrule
\end{tabular}
\end{block}
\end{column}

\end{columns}
\end{frame}

% --- B2: Seed asymmetry ---
\begin{frame}{Backup B2: Asimetría de semillas y caching}
\small
\textbf{Observación:} 80\% de combinaciones LLM $\times$ clasificador determinista
tienen \textbf{varianza cero} entre semillas (la generación está cacheada).

\vspace{0.4em}
\begin{columns}[T]

\begin{column}{0.55\textwidth}
\textbf{Validación rigurosa:}
\begin{itemize}\setlength\itemsep{0pt}
  \item Filtramos a clasificadores estocásticos\\(MLP, RandomForest).
  \item $\Delta$F1 = $+1{.}79$ pp.
  \item $p = 0{.}0015$, $d = 0{.}52$.
  \item Significancia se mantiene.
\end{itemize}
\end{column}

\begin{column}{0.45\textwidth}
\setbeamercolor{block title}{fg=white,bg=uaiamber}
\begin{block}{Implicancia}
\scriptsize
Los CIs subestiman la varianza real, pero las pruebas $t$ pareadas \textbf{siguen siendo válidas}.
\end{block}
\end{column}

\end{columns}
\end{frame}

% --- B3: Closed-loop (iterative prompting) ---
\begin{frame}{Backup B3: Closed-loop vs.\ pasada única}
\small
\textbf{Hipótesis exploratoria:} ¿mejora si re-prompteamos al LLM con feedback de los rechazos?

\vspace{0.5em}
\begin{tabular}{@{}lrr@{}}
\toprule
Estrategia                          & $\Delta$F1 vs.\ SMOTE & Win-rate \\
\midrule
Closed-loop (5 iteraciones)         & $+1{.}12$             & 52\% \\
Pasada única \alert{$\star$}        & $+2{.}08$             & 80\% \\
\bottomrule
\end{tabular}

\vspace{0.6em}
\textbf{¿Por qué falla closed-loop?}
\begin{itemize}\setlength\itemsep{0pt}
  \item Tasas de aceptación \textbf{decrecen} con iteraciones (-12.3 pp).
  \item Fallo dominante: \texttt{DISTANCE\_OUTLIER} (77\%).
  \item Cuello de botella: \textbf{el LLM}, no el prompt.
\end{itemize}

\vspace{0.4em}
\begin{insight}
\scriptsize Pasada única + filtro simple $>$ ingeniería iterativa de prompt.
\end{insight}
\end{frame}

% --- B4: Geometric vs Teacher ---
\begin{frame}{Backup B4: Filtrado geométrico vs.\ \textit{teacher}}
\small
\textbf{Comparación con un filtrado supervisado (teacher classifier):}

\vspace{0.5em}
\begin{tabular}{@{}lrrr@{}}
\toprule
Filtro & $\Delta$F1 vs.\ SMOTE & $p$ & Costo extra \\
\midrule
Geométrico (binary) & \alert{$+2{.}78$} & $<0{.}001$ & ninguno \\
Teacher (binary)    & $+2{.}43$         & $<0{.}001$ & entrenar clf adicional \\
\bottomrule
\end{tabular}

\vspace{0.6em}
\begin{insight}
\scriptsize El filtro geométrico es \textbf{más general} (no se especializa en separabilidad lineal) y \textbf{más eficiente} (no requiere entrenamiento extra).
\end{insight}
\end{frame}

% --- B5: Tabla completa de baselines modernos ---
\begin{frame}{Backup B5: Comparación con baselines modernos}
\small
\textbf{1.890 configuraciones, 10 métodos, 3 clasificadores:}

\vspace{0.5em}
\begin{tabular}{@{}lrrr@{}}
\toprule
Método                         & $\Delta$F1 vs.\ SMOTE & $p$ Bonf. & Win-rate \\
\midrule
\textbf{Soft-weighted (ours)}  & \alert{$+2{.}61$}     & $<0{.}0001$ & 88.9\% \\
\textbf{Binary-filter (ours)}  & $+2{.}39$             & $<0{.}0001$ & 79.4\% \\
No augmentation                & $+0{.}64$             & $<0{.}05$   & 56.0\% \\
Embedding mixup                & $-0{.}01$             & n.s.        & 51.0\% \\
Contextual BERT                & $-0{.}26$             & n.s.        & 48.0\% \\
T5 paraphrase                  & $-0{.}35$             & n.s.        & 47.0\% \\
Random oversample              & $-0{.}32$             & n.s.        & 47.0\% \\
EDA                            & $-0{.}49$             & $<0{.}05$   & 41.0\% \\
Back-translation               & $-0{.}50$             & $<0{.}05$   & 40.0\% \\
\bottomrule
\end{tabular}

\vspace{0.5em}
\begin{insight}
\scriptsize Todas las técnicas modernas \textbf{pierden} contra SMOTE; solo el filtrado geométrico gana.
\end{insight}
\end{frame}


\end{document}
